% =============================================================
%  3 仿真结果及分析
% =============================================================
\section{仿真结果及分析}

\subsection{仿真环境及参数}

将环境设置为使用 Python 3.7、PyCharm 2022 和 TensorFlow 1.14。默认参数设置如
表~\ref{tab:params} 所示，除非另有说明，否则在以下仿真中将保持不变。

\begin{table}[H]
  \centering
  \caption{参数说明和默认值设置}
  \label{tab:params}
  \renewcommand{\arraystretch}{1.2}
  \begin{tabularx}{\linewidth}{X c}
    \toprule
    \multicolumn{1}{c}{参数} & \multicolumn{1}{c}{数值} \\
    \midrule
    IoT 设备数量 & 10 \\
    基站数量 & 3 \\
    每回合时隙数 & 100 \\
    最大输入数据量/Mbit & 2 \\
    最小输入数据量/Mbit & 5 \\
    最大计算复杂度/(CPU cycles$\cdot$bit$^{-1}$) & 100 \\
    最小计算复杂度/(CPU cycles$\cdot$bit$^{-1}$) & 1\,000 \\
    最大传输速率/(Gbit$\cdot$s$^{-1}$) & 1 \\
    最小传输速率/(Gbit$\cdot$s$^{-1}$) & 2 \\
    任务处理的最大容忍时延/s & 10 \\
    每个时隙的长度/s & 15 \\
    最大块大小/Mbit & 0.5 \\
    最小块大小/Mbit & 1 \\
    传播时延/s & 0.01 \\
    \bottomrule
  \end{tabularx}
\end{table}

为了验证区块链驱动的边缘计算系统联合优化算法的性能，首先，将所提算法在不同学习率下
对评论家网络和演员网络的收敛性进行分析。其次，与随机接入算法和 AC\textsuperscript{[27]}
算法在任务处理的奖励和能量消耗等方面进行比较。

\subsection{结果分析}

考虑评论家网络的 3 种不同的学习率，将演员网络的学习率设置为 $\eta^{\mu}$。同时，
保持评论家网络的学习率，并分析 3 种不同学习率对演员网络的影响。所提算法在不同学习率下
评论家网络和演员网络的收敛性具体如图~\ref{fig:convergence-critic}~图~\ref{fig:convergence-actor}
所示。

\begin{figure}[H]
  \centering
  \fbox{\parbox{0.9\linewidth}{\centering
    \vspace{1.8cm} 图~\ref{fig:convergence-critic} 占位：评论家网络收敛曲线 \vspace{1.8cm}}}
  \caption{所提算法在不同学习率下对评论家网络的收敛性}
  \label{fig:convergence-critic}
\end{figure}

\begin{figure}[H]
  \centering
  \fbox{\parbox{0.9\linewidth}{\centering
    \vspace{1.8cm} 图~\ref{fig:convergence-actor} 占位：演员网络收敛曲线 \vspace{1.8cm}}}
  \caption{所提算法在不同学习率下对演员网络的收敛性}
  \label{fig:convergence-actor}
\end{figure}

由图~\ref{fig:convergence-critic} 和图~\ref{fig:convergence-actor} 可以看出，这些算法
收敛曲线分别在第 200 个回合附近收敛，表明在解决复杂问题时，所提算法表现出了快速的收敛特性。

将所提算法和随机接入算法及 AC 算法在任务处理奖励和能量消耗等方面进行比较分析。
为了最小化能量消耗，通过总能量消耗的倒数确定奖励的大小，因此较高的奖励值表示系统性能较好，
在降低能耗方面性能更优越。3 种算法的综合性能对比结果如图~\ref{fig:comparison} 所示。

\begin{figure}[H]
  \centering
  \fbox{\parbox{0.9\linewidth}{\centering
    \vspace{1.8cm} 图~\ref{fig:comparison} 占位：3 种算法综合性能对比 \vspace{1.8cm}}}
  \caption{3 种算法的综合性能对比结果}
  \label{fig:comparison}
\end{figure}

由图~\ref{fig:comparison} 可以看出，当所有参数都采用默认设置时，所提算法在收敛性上更稳定，
且能实现最大的收益。AC 算法和随机接入算法虽然能够收敛，但其收益明显低于所提算法，且稳定性较差。
因此，所提算法在收敛后，其奖励几乎保持稳定，具有强大的鲁棒性和稳定性。
